% CLAUDE REPORT::BRIDGE-PINCER-REFLECTION-0001::2026-11-14
\documentclass[11pt]{article}
\usepackage[margin=1in]{geometry}
\usepackage{amsmath,amssymb,amsthm}
\usepackage{booktabs}
\usepackage[hidelinks]{hyperref}
\newtheorem{lemma}{Lemma}
\newtheorem{proposition}[lemma]{Proposition}
\newtheorem{corollary}[lemma]{Corollary}
\newcommand{\SAT}{\mathrm{SAT}}
\newcommand{\poly}{\operatorname{poly}}
\title{The Bridge--Lower-Bound Pincer and the Reflection Ladder for $\overline W_R(N)$}
\author{Claude scheduled worst-case-derivation run (supporting all five roles)}
\date{14 November 2026}
\begin{document}\maketitle

\section*{Status}
\textbf{P versus NP remains OPEN.} The lemmas below are PROVED/EXACT in the stated models
(most are short consequences of KNOWN THEOREMS). The load-bearing bridge B3 for any fixed
system whose lower bound is not already known stays \textbf{OPEN}. Nothing here is an
unrestricted SAT lower bound. The report writes no role-owned region of root \texttt{main.tex}:
this scheduled run is not one of LSPV-1..5, and it adds evidence without changing any task,
schedule, role, or allocation.

\section*{Shared objective (unchanged)}
For a fixed sound, complete, polynomial-time-checkable calculus $R$ with cost $\mathrm{cost}_R$,
\[
w_R(F)=\min_{\pi:F\vdash_R\bot}\mathrm{cost}_R(\pi),\qquad
\overline W_R(N)=\max_{F\in\mathrm{UNSAT},\,|\langle F\rangle|\le N} w_R(F).
\]
$\mathrm{LB}_R$: $\overline W_R(N)\ne N^{O(1)}$. \quad $\mathrm{B3}_R$: $\SAT\in P\Rightarrow\overline W_R(N)\le\poly(N)$.
Choosing the optimal derivation for free (the minimum over $\pi$) is a valid lower bound on the work of any
procedure that \emph{outputs or executes} an $R$-derivation. The open gap is a decider that never
produces an $R$-derivation. B3 is exactly the statement that closes that gap.

\section*{1. The pincer (PROVED/EXACT)}
\begin{lemma}[Pincer]\label{pincer}
(a) For every fixed $R$: if $\mathrm{LB}_R$ holds, then $\mathrm{B3}_R\iff\SAT\notin P$.
(b) $\exists R\,[\mathrm{B3}_R\wedge \mathrm{LB}_R]\iff \SAT\notin P$.
\end{lemma}
\begin{proof}
(a) If $\mathrm{LB}_R$ holds, the consequent of $\mathrm{B3}_R$ is false, so $\mathrm{B3}_R$ is equivalent to the negation of its antecedent.
(b) ($\Rightarrow$) Modus tollens. ($\Leftarrow$) Take $R=$ Resolution with node cost. Haken (1985) gives $\mathrm{LB}_R$, and $\mathrm{B3}_R$ is vacuously true when $\SAT\notin P$.
\end{proof}
\noindent\emph{Consequence.} The strategy proves something only if $\mathrm{B3}_R$ is proved
\emph{constructively}, by turning a hypothetical polynomial decider into polynomial $R$-derivations, for an $R$ whose
$\mathrm{LB}_R$ is \emph{not} already known. For Resolution, Cutting Planes, Polynomial Calculus,
SOS, bounded-depth Frege and $\mathrm{AC}^0[p]$-Frege (i.o.\ lower bounds, Lu--Santhanam--Tzameret, ECCC TR25-134), B3 is just
P$\ne$NP restated. The corpus already records this for Resolution and padded systems (LSPV-1/335 line, LSPV-5/409). Part (b) adds the general existential form.

\begin{lemma}[Monotonicity]\label{mono}
Suppose $R'$ polynomially simulates $R$ with respect to cost: there are a polynomial-time map $\pi\mapsto\pi'$
and a polynomial $q$ with $\mathrm{cost}_{R'}(\pi')\le q(\mathrm{cost}_R(\pi)+|F|)$. Then
$\mathrm{B3}_R\Rightarrow\mathrm{B3}_{R'}$ and $\mathrm{LB}_{R'}\Rightarrow\mathrm{LB}_R$.
\end{lemma}
\begin{proof}
We have $w_{R'}(F)\le q(w_R(F)+|F|)$. Take the maximum over $F$ to get $\overline W_{R'}(N)\le q(\overline W_R(N)+N)$ (the maximum of a monotone $q$ is attained at the same $F$). Both implications follow.
\end{proof}
\noindent So moving $R$ up the simulation order makes the bridge easier and the lower bound harder, and
by Lemma~\ref{pincer}(b) no position is unconditionally better than another. The useful window is
$R\succeq\mathrm{EF}$, where a constructive B3 route exists (LSPV-1/357 and /592). Lower bounds are known only
far below that, up to $\mathrm{AC}^0[p]$-Frege. Today the window
$\{R:\ \text{constructive B3}\}\cap\{R:\ \text{LB known}\}$ is empty.

\section*{2. Reflection ladder: the exact weakest constructive bridge}
Fix a sound r.e.\ theory $T\supseteq\mathrm{PV}$ (for example $S^1_2$, PA, ZFC) and its propositional system
$P_T$: a proof of UNSAT $F$ is a $T$-proof of $\forall y\,\neg\mathrm{Eval}(\ulcorner F\urcorner,y)$.
Cost is encoded length. For a polynomial-time $A$, let $\mathrm{RS}_A$ be its rejection soundness, as in LSPV-1/357.

\begin{proposition}[Provable rejecter $\Rightarrow$ $P_T$ bounded]\label{up}
If some polynomial-time $A$ satisfies $T\vdash\mathrm{RS}_A$, then $\overline W_{P_T}(N)\le\poly(N)$.
\end{proposition}
\begin{proof}
Concatenate three parts: the fixed constant-size $T$-proof of $\mathrm{RS}_A$; its instantiation at $\ulcorner F\urcorner$; and a
$T$-proof of the true bounded statement $\mathrm{Reject}_A(\ulcorner F\urcorner)$. The last part has length $\poly(\mathrm{time}_A(F))$, because
PV proves correct polynomial-time computations with proofs polynomial in the length of the trace. Modus ponens gives the result.
\end{proof}

\begin{proposition}[Levin closure, one reflection step down]\label{down}
Assume $\SAT\in P$ and $\overline W_{P_T}(N)\le p(N)$. Then there is a polynomial-time $A^\ast$ with
$T+\mathrm{Rfn}_{P_T}\vdash\mathrm{RS}_{A^\ast}$, where $\mathrm{Rfn}_{P_T}$ is the local reflection principle for $P_T$.
\end{proposition}
\begin{proof}
"Is there a $P_T$-proof of $F$ of length $\le p(|F|)$ extending a given prefix?" is an NP question. With $\SAT\in P$ and self-reducibility, a polynomial-time $A^\ast$ either finds such a proof and checks it, or declares that none exists. $A^\ast$ rejects only after checking a proof. So $\mathrm{RS}_{A^\ast}$ follows from $\mathrm{Rfn}_{P_T}$ together with PV-provable correctness of the checker.
\end{proof}

\begin{corollary}[Ladder form of B3]\label{ladder}
Let $T_0=T$ and $T_{k+1}=T_k+\mathrm{Rfn}_{P_{T_k}}$. Under $\SAT\in P$, the condition
"$P_{T_k}$ is polynomially bounded for some $k$" holds if and only if "some polynomial-time $A$ has
$T_k\vdash\mathrm{RS}_A$ for some $k$". Hence, for strong fixed systems, the weakest constructive bridge is
\[
\boxed{\SAT\in P\ \Longrightarrow\ \exists A\in P\ \exists k:\ T_k\vdash\mathrm{RS}_A}
\qquad(\text{provable-collapse bridge, OPEN}).
\]
\end{corollary}
\noindent The Gödel firewall (LSPV-1/357) shows that \emph{some} correct polynomial deciders are not
$T$-certifiable. It does not show that \emph{all} of them fail. The boxed statement is therefore not refuted.
Lemma~\ref{pincer} does not apply here, because $\mathrm{LB}_{P_T}$ is wide open. The boxed statement
is strictly weaker than the PV/EF route of /357: it may use any $T_k$. The price is that $\mathrm{LB}$ is then
needed for $P_{T_k}$, a system at least as strong as ZFC-proofs if $T=\mathrm{ZFC}$.
Pich--Santhanam (FOCS 2019; arXiv 2312.08163) and Khaniki--Pich--Sokolov (CCC 2026) give
evidence that superpolynomial lower bounds for EF and stronger systems face inherent provability obstacles.
\emph{Classification:} the ladder is EXACT within its hypotheses. Its use as a P$\ne$NP route is CONDITIONAL on
the boxed bridge and on $\mathrm{LB}_{P_{T_k}}$, and both are OPEN.

\section*{3. Hard-instance survival (KNOWN THEOREM summary)}
\begin{tabular}{@{}lll@{}}\toprule
Family & Killed by & Survives \\\midrule
PHP$^{n+1}_n$ & Cutting Planes (poly), EF (poly, Cook--Reckhow) & Res, PC (Razborov), bd-depth Frege \\
Tseitin, odd charge & Gaussian elimination / $\mathbb F_2$-PC, Res(lin) & Res (Urquhart), PC over char $\ne 2$ \\
random 3-CNF, $m=\Delta n$ & --- (no known poly in any studied system) & Res, PC, SOS (Grigoriev), CP (Fleming et al.) \\
Clique-coloring & EF-style circuit reasoning & CP (Pudlák 1997, interpolation) \\
\bottomrule\end{tabular}

\noindent Algebraic reasoning, circuit compression, shared DAGs and adaptive switching all amount to moving to a
stronger $R$, which is covered by Lemma~\ref{mono}. Random 3-CNF is the only standard family not known to fall to any of
them. Its EF hardness is a conjecture (Krajíček's proof-complexity-generator program), so it is CANDIDATE only.

\section*{CORPUS REVIEWED}
AGENTS.md (including the knot/link and world-ledger notes), the 2026-11-02 worst-case minimum-derivation refocus,
LSPV-1/335, /357 (Gödel firewall), /395 (noncomputable profile), /540, /592 (ER strength barrier),
LSPV-5/409 (identity-system boundary), Codex 2026-11-11 SAT-polynomial audit and transfer report,
the world ledger (rows L13, L34), and the most recent TMR-1..5 commit log (MCSP track).
\section*{DUPLICATION / CONTRADICTION AUDIT}
Already in the corpus: Haken vacuity for Resolution B3, the PV/EF sufficiency route (/357), and the ER simulation (/592).
New in this report: the existential pincer (Lemma 1b), cost-monotonicity as an explicit tradeoff, the Levin closure,
and the reflection-ladder equivalence (Prop.~4, Cor.~5). No contradiction with root \texttt{main.tex} statuses was found.
\section*{WHAT CHANGED RELATIVE TO ROOT MAIN}
No edit (this run has no ownership region). Suggested for LSPV-1/LSPV-5: record Cor.~5 as the exact
weakest constructive form of B3 for strong fixed systems.
\section*{ADOPTED / REJECTED / RETIRED PRIOR COMPONENTS}
Adopted: /357, /592, /409. Retired as a route: any B3 for a system with known $\mathrm{LB}$; it is a reformulation (Lemma 1).
\section*{STALE OR OVERSTATED CLAIMS FOUND}
None found in the files reviewed.
\section*{Barrier transfer and PQ ledger}
Control theory: $\mathrm{RS}_A$ plays the role of a barrier certificate. The SOS-certificate degree lower bounds are a
\emph{directly transferable} instance of Lemma 1(a): B3 for SOS is equivalent to P$\ne$NP. PQ item 6 (special soundness):
analogy only. "Reject only after a checked proof" (Prop.~4) is the extraction pattern, and it carries no
hardness assumption. Other PQ items: FAILED TRANSFER, not relevant here.
\section*{Next obligations}
(1) LSPV-1: decide whether the boxed bridge relativizes. It cannot be a purely relativizing argument, since there are oracles with $P=NP$.
(2) LSPV-2 and LSPV-3: test random 3-CNF in the native calculus and in Res(lin).
(3) LSPV-4: charge the trace length in Prop.~3 in the matched machine model.
(4) LSPV-5: red-team Prop.~4 on the self-reducibility step for malformed inputs.

\medskip\noindent\textbf{P versus NP remains OPEN.}
\end{document}
% END_OF_VERIFIED_REPORT::CLAUDE::BRIDGE-PINCER-REFLECTION-0001
